\subsection{融合误差与最优融合权重}
\label{sec:error_model}

本节将每个邻居视为独立的观测来源，引入测量与对齐误差，刻画各邻居``应该贡献多少''。由式~\eqref{eq:final_shift}，定义最终对齐误差
\begin{equation}
  \epsilon_{pq}\triangleq\hat k_{pq}-(\mu_p-\mu_q)
  =(e_p-e_q)+r_{pq}+\hat\delta_{pq},
  \label{eq:epsilon}
\end{equation}
其中 $r_{pq}=k_{pq}^{M}-\Delta_{pq}\in[-\tfrac12,\tfrac12]$ 为取整残差。对齐后的邻居回波中心位于 $\mu_q+\hat k_{pq}=\mu_p+\epsilon_{pq}$，其观测为
\begin{equation}
  y_{pq}=\mu_p+\epsilon_{pq}+n_{pq},
  \label{eq:neighbor_obs}
\end{equation}
其中 $n_{pq}$ 为Coates校正后的直方图上回波中心的估计误差，方差由式~\eqref{eq:coates_var}经误差传播给出，记为 $v_q$。对齐误差的方差记为 $\sigma_{\epsilon,pq}^2$。当二者条件独立且对齐偏差已校准时，单个邻居的总不确定度为
\begin{equation}
  v_{pq}^{\mathrm{eff}}=v_q+\sigma_{\epsilon,pq}^{2}.
  \label{eq:veff}
\end{equation}
若各邻居的总不确定度逐对已知，则最优线性融合为逆方差加权，邻居 $q$ 的权重为
\begin{equation}
  c_{pq}=\frac{1}{v_{pq}^{\mathrm{eff}}}
  =\frac{1}{v_q+\sigma_{\epsilon,pq}^{2}}.
  \label{eq:c_pq}
\end{equation}
然而这一前提无法满足，$\sigma_{\epsilon,pq}^{2}$ 的逐对取值依赖未知的真实位移与匹配细节，无法从单次观测直接计算；即便由配对数据整体估计，得到的也只是平均散布，而非``这一对''邻居的对齐可靠性。平均散布至多可充当权重中的基底项，逐对的偏离仍需另行处理。
